\section{Evaluation}\label{sec:evaluation}

\begin{table*}[t]
\centering
\caption{Intrusion detection performance across three graph-based security datasets. The supervised GNN uses labeled attacks during training. GraphCL and \sysname{} are benign-only representation learning methods.}
\label{tab:main_results}
\resizebox{\textwidth}{!}{%
\begin{tabular}{llcccc}
\toprule
Dataset & Model & Accuracy & Precision & Recall & F1 \\
\midrule
StreamSpot & Supervised GNN & $0.989 \pm 0.016$ & $1.000 \pm 0.000$ & $0.933 \pm 0.094$ & $0.963 \pm 0.052$ \\
StreamSpot & GraphCL & $0.750 \pm 0.036$ & $0.353 \pm 0.054$ & $0.733 \pm 0.377$ & $0.450 \pm 0.154$ \\
StreamSpot & \sysname{} & $0.917 \pm 0.023$ & $0.862 \pm 0.037$ & $0.996 \pm 0.006$ & $0.923 \pm 0.019$ \\
\midrule
GraSec-IoT & Supervised GNN & $0.988 \pm 0.003$ & $0.997 \pm 0.002$ & $0.990 \pm 0.004$ & $0.994 \pm 0.002$ \\
GraSec-IoT & GraphCL & $0.953 \pm 0.006$ & $0.970 \pm 0.006$ & $0.980 \pm 0.008$ & $0.975 \pm 0.003$ \\
GraSec-IoT & \sysname{} & $0.742 \pm 0.124$ & $0.718 \pm 0.129$ & $0.909 \pm 0.064$ & $0.790 \pm 0.065$ \\
\midrule
Unicorn Wget & Supervised GNN & $0.867 \pm 0.060$ & $0.600 \pm 0.490$ & $0.267 \pm 0.226$ & $0.367 \pm 0.306$ \\
Unicorn Wget & GraphCL & $0.560 \pm 0.200$ & $0.123 \pm 0.101$ & $0.600 \pm 0.490$ & $0.203 \pm 0.167$ \\
Unicorn Wget & \sysname{} & $0.515 \pm 0.062$ & $0.537 \pm 0.071$ & $0.710 \pm 0.215$ & $0.583 \pm 0.053$ \\
\bottomrule
\end{tabular}%
}
\end{table*}

This section evaluates whether \sysname{} provides effective graph-level intrusion detection across datasets, baselines, and label-availability settings.

\subsection{Research Questions}

We evaluate three research questions:

\begin{itemize}
    \item RQ1: How does \sysname{} compare with supervised and benign-only graph baselines across intrusion detection datasets?

    \item RQ2: Does topology-aware contrastive scoring improve over standard graph contrastive learning?

    \item RQ3: How do the methods balance attack recall and alert precision?
\end{itemize}

\subsection{Datasets}

We evaluate graph-level intrusion detection on three datasets. Each dataset is converted into a set of graphs. Each graph has a benign or malicious label for validation and testing. During representation learning, benign-only methods use only benign training graphs.

\begin{itemize}
    \item \textbf{StreamSpot:} contains heterogeneous system-activity provenance graphs. These graphs represent relationships among system entities and events. StreamSpot is useful for testing whether a detector can identify malicious behavior in structured system activity \cite{manzoor2016streamspot}.

    \item \textbf{GraSec-IoT:} contains graph snapshots from an Internet of Things security setting. It includes communication behavior and attack activity in an IoT environment \cite{hamouche2024graseciot}. This dataset tests whether the method works on IoT-style communication graphs.

    \item \textbf{Unicorn Wget:} is a graph-based security dataset built from system activity collected from a Wget attack scenario. It provides graph samples rather than tabular flow records, so it fits our graph-level detection setting \cite{unicornwget}.
\end{itemize}

For all datasets, we use the same split logic. Benign-only methods train their encoders using only benign training graphs. The validation set contains benign and malicious graphs and is used only to choose the anomaly threshold. The test set is used once for final evaluation. We use the same splits and random seeds for all methods.

\subsection{Baselines}

We compare \sysname{} with one label-rich reference model and one benign-only representation learning baseline.

\begin{itemize}
    \item \textbf{Supervised GNN.} This supervised graph neural network is trained with labeled benign and malicious graphs. We treat it as a label-rich reference baseline. It is not directly comparable to benign-only representation learning methods because it uses attack labels during training, but it shows what performance is possible when labels are available.

    \item \textbf{GraphCL.} This benign-only graph contrastive learning baseline learns graph embeddings from augmented graph views without using attack labels during representation learning \cite{you2020graphcl}. It tests whether standard graph contrastive learning is sufficient for this task.
\end{itemize}

\subsection{Experimental Setup}

All methods use the same split logic for a given dataset and seed. StreamSpot and GraSec-IoT results aggregate seeds 42, 43, and 44; Unicorn Wget results aggregate seeds 42 through 46. The threshold is selected separately on the validation set for each benign-only run before being applied once to the corresponding test set.

The encoder is trained only on benign training graphs. Malicious graphs are used only for validation threshold selection and final testing. This prevents attack labels from being used during representation learning.

For \sysname{}, we use 12 training epochs, hidden dimension 64, embedding dimension 64, temperature 0.2, maximum 512 nodes per graph, node drop rate 0.10, edge drop rate 0.05, feature mask rate 0.10, topology mask rate 0.10, topology noise 0.10, and $k=5$ nearest neighbors for anomaly scoring. For the joint graph-topology model, we use $\alpha=0.5$ and $\beta=0.5$ to weight the graph and topology contrastive losses. Baseline hyperparameters follow their saved experimental configurations rather than being assumed identical to \sysname{}.

\subsection{Metrics}

We report accuracy, precision, recall, and F1 score. Let TP be true positives, TN be true negatives, FP be false positives, and FN be false negatives. Precision is $\frac{TP}{TP+FP}$ and measures how many predicted attacks are actually attacks. Recall is $\frac{TP}{TP+FN}$ and measures how many attacks are detected. F1 score is the harmonic mean of precision and recall.

F1 is useful because intrusion detection needs both high precision and high recall. Precision is important because too many false alerts can make an intrusion detection system difficult to use in practice. Recall is important because missed attacks can leave a system exposed.

\subsection{Results by Research Question}

Table~\ref{tab:main_results} reports intrusion detection performance across StreamSpot, GraSec-IoT, and Unicorn Wget. The supervised GNN uses labeled benign and malicious graphs during training. GraphCL and \sysname{} train their encoders using only benign graphs and use labeled validation data only for threshold selection.

\noindent\textbf{RQ1: Comparison with supervised and benign-only baselines.}
The supervised GNN achieves the highest F1 on StreamSpot ($0.963 \pm 0.052$) and GraSec-IoT ($0.994 \pm 0.002$), while \sysname{} achieves the highest F1 on Unicorn Wget ($0.583 \pm 0.053$ versus $0.367 \pm 0.306$ for the supervised GNN). Among the benign-only methods, \sysname{} has higher F1 on StreamSpot ($0.923 \pm 0.019$ versus $0.450 \pm 0.154$) and Unicorn Wget ($0.583 \pm 0.053$ versus $0.203 \pm 0.167$), whereas GraphCL leads on GraSec-IoT ($0.975 \pm 0.003$ versus $0.790 \pm 0.065$). Thus, \sysname{} is competitive in two benign-only comparisons but is neither uniformly superior to GraphCL nor a replacement for label-rich learning when representative attack labels are available.

\noindent\textbf{RQ2: Effect of topology-aware contrastive scoring.}
Topology-aware contrastive scoring does not improve over standard GraphCL uniformly. \sysname{} has higher F1 by $0.473$ on StreamSpot and $0.380$ on Unicorn Wget, but lower F1 by $0.185$ on GraSec-IoT. This heterogeneity indicates that the value of degree- and closeness-based topological summaries depends on the dataset and graph construction. A plausible explanation is that topology adds complementary information when anomalous behavior changes connectivity at graph scale, while the standard graph representation is already sufficient for the GraSec-IoT split. Because the methods also differ in architecture and training configuration, these comparisons do not isolate topology as the sole causal factor.

\noindent\textbf{RQ3: Precision--recall balance.}
On StreamSpot, \sysname{} combines recall of $0.996 \pm 0.006$ with precision of $0.862 \pm 0.037$, whereas GraphCL has lower and much more variable recall ($0.733 \pm 0.377$) and precision ($0.353 \pm 0.054$). On GraSec-IoT, GraphCL provides the stronger balance, with precision of $0.970 \pm 0.006$ and recall of $0.980 \pm 0.008$; \sysname{} retains high recall ($0.909 \pm 0.064$) but lower precision ($0.718 \pm 0.129$). On Unicorn Wget, \sysname{} improves the benign-only balance, but both benign-only methods and the supervised GNN exhibit substantial weaknesses or variability. For resilient monitoring, the preferred operating point remains deployment-dependent: missed attacks favor recall, whereas analyst capacity and alert fatigue favor precision.

Overall, RQ1 shows that the supervised reference has the highest F1 on StreamSpot and GraSec-IoT, while \sysname{} leads on Unicorn Wget and is the stronger benign-only method on StreamSpot and Unicorn Wget. RQ2 shows that topology-aware scoring is dataset-dependent rather than uniformly beneficial. RQ3 shows a strong TRACE recall--precision balance on StreamSpot, a stronger GraphCL balance on GraSec-IoT, and challenging, variable behavior for all methods on Unicorn Wget.
